%!TEX root = ../../main.tex

\chapter{Handlungsempfehlungen für die Praxis}
\label{ch:handlungsempfehlungen}

Die in Abschnitt~\ref{sec:beantwortung-ff} gegebenen Antworten auf die Forschungsfragen lassen sich unmittelbar in Empfehlungen für die Gestaltung \ac{KI}-gestützter \ac{ETL}-Strecken überführen. Dieses Kapitel leitet sie aus den Befunden ab und ordnet sie den beiden Fragen zu. Sie werden entlang dreier Entscheidungen entwickelt, die in einem Vorhaben nacheinander zu treffen sind: ob ein \ac{LLM} für eine gegebene Aufgabe überhaupt das geeignete Mittel ist (Abschnitt~\ref{sec:eignung}), wie Modell, Ausführungsmodus und Prompt zu wählen sind (Abschnitt~\ref{sec:empfehlungen-auswahl}) und wie sich beide Verfahrensarten in einer Gesamtarchitektur verbinden lassen (Abschnitt~\ref{sec:hybride-architekturen}). Für alle gilt vorab dieselbe Einschränkung, die Abschnitt~\ref{sec:limitationen} ausführt: Die Empfehlungen sind Ausgangspunkte für eine eigene Erprobung. Weder die konkreten Zahlenwerte noch die vermuteten Wirkungszusammenhänge sind ohne Prüfung auf andere Daten und Systeme übertragbar.

\section{Eignung von LLMs für unterschiedliche ETL-Szenarien}
\label{sec:eignung}

Der über alle Untersuchungsebenen hinweg stabilste Befund dieser Arbeit ist, dass die Eignung eines \ac{LLM} weit stärker von der Art der Aufgabe abhängt als von der Wahl des Modells. Zwischen dem besten und dem schwächsten Cloud-Modell liegen rund acht Prozentpunkte, zwischen der am besten und der am schlechtesten gelösten Aufgabe hingegen nahezu hundert. Die erste Frage in einem \ac{ETL}-Vorhaben lautet daher nicht, welches Modell einzusetzen ist, sondern für welche Teilaufgaben sich der Einsatz überhaupt lohnt \autocite{zhaoSurveyLargeLanguage2026,2026AIIndex}.

Für diese Entscheidung hat sich in den Ergebnissen eine einzige Unterscheidung als tragfähig erwiesen: ob die Verarbeitungsvorschrift vollständig vorab formulierbar ist oder nicht. Aufgaben der ersten Art (das Entfernen von Leerzeichen, das Vereinheitlichen bekannter Datumsformate, das Entfernen exakter Duplikate, die schlüsselbasierte Entdoppelung, Typkonvertierungen, das Berechnen abgeleiteter Spalten) werden von einer regelbasierten Implementierung definitionsgemäß korrekt gelöst, deterministisch, in Sekundenbruchteilen und ohne laufende Kosten. Ein \ac{LLM} erreicht hier bestenfalls dasselbe Ergebnis und bringt zusätzlich Nichtdeterminismus, Kosten und Prüfaufwand mit sich. Sein Beitrag liegt bei diesen Aufgaben ausschließlich in der \emph{Erstellung} der Regel, nicht in ihrer Ausführung, ein Unterschied, der in der Diskussion über \ac{KI}-gestützte Datenverarbeitung häufig verwischt wird \autocite{coteDataCleaningMachine2024,aAIEnhancedETLProcesses2025}.

Aufgaben der zweiten Art sind dadurch gekennzeichnet, dass die Vorschrift eine Aufzählung erfordern würde, die niemand vollständig angeben kann. Die Vereinheitlichung der Länderangaben ist der klarste Fall: Eine Zuordnungstabelle erfasst nur, was zuvor jemand eingetragen hat. Die regelbasierte Vergleichspipeline scheiterte an genau drei Schreibweisen, nämlich \enquote{Holland}, \enquote{England} und dem Tippfehler \texttt{Deutshcland}. Genau hier liegt der originäre Beitrag des Modells, das erschließt, dass \enquote{Holland} umgangssprachlich für die Niederlande steht, ohne dass diese Gleichsetzung je hinterlegt wurde. Ähnlich, wenn auch weniger ausgeprägt, liegt der Fall bei der unscharfen Duplikaterkennung, wo die Code-Generierung im besten Lauf die regelbasierte Lösung übertrifft: Die Entscheidung über zwei konkrete Datensätze ist einem Modell unmittelbar zugänglich, während ihre Verallgemeinerung zu einem Schwellenwert Information vernichtet.

Diese Zuordnung bleibt jedoch unvollständig, solange sie den Ausführungsmodus außer Acht lässt, denn der Vorteil bei semantischen Aufgaben tritt nur dann ein, wenn das Modell sein Wissen auf den einzelnen Wert anwenden kann. In generierten Code gefasst, fiel dieselbe Aufgabe auf das Niveau der regelbasierten Lösung zurück, weil die erzeugten Skripte ebenfalls unvollständige Zuordnungstabellen verwendeten. Andere Formen generierter Programme sind damit nicht ausgeschlossen. Die Eignungsfrage ist deshalb nicht mit \enquote{\ac{LLM} oder Regel} beantwortet, sondern erst mit der Angabe des Modus. Tabelle~\ref{tab:eignungsmatrix} fasst sie zusammen.

\begin{table}[htbp]
    \centering
    \small
    \caption[Eignungsmatrix nach Aufgabentyp]{Empfohlenes Verfahren je Aufgabentyp auf Grundlage der Ergebnisse aus Kapitel~\ref{ch:ergebnisse}. Die letzte Spalte nennt den jeweils tragenden Messwert.}
    \label{tab:eignungsmatrix}
    \begin{tabularx}{\textwidth}{@{}p{2.4cm}XXX@{}}
        \toprule
        \textbf{Aufgabentyp} & \textbf{Merkmal} & \textbf{Empfehlung} & \textbf{Beleg} \\
        \midrule
        Syntaktisch regelbar & Vorschrift vollständig aufzählbar & Regelbasiert ausführen, \ac{LLM} allenfalls zur Erstellung des Codes & Regelbasiert 100\,\%, bestes \ac{LLM} ebenfalls 100\,\% \\
        \addlinespace
        Semantisch, wertweise & Entscheidung je Einzelwert, Verallgemeinerung verlustbehaftet & \ac{LLM} im Direkt-Modus, gezielt als einzelner Schritt & Ländercodes: Direkt 97,1--100\,\% gegen Code-Generierung 94,8\,\% \\
        \addlinespace
        Rechnerisch aggregierend & viele Datensätze, exakte Arithmetik & Regelbasiert oder \ac{LLM} im Code-Generierungs-Modus, Direkt-Modus ungeeignet & Abgeleitete Spalten: Code-Generierung 100\,\%, Direkt-Modus am Ausgabelimit abgebrochen \\
        \addlinespace
        Zusammen-  gesetzt & mehrere Entscheidungen in einem Schritt & In Einzelschritte zerlegen, dann je Schritt entscheiden & Einzelaufgabe 0,1\,\%, zerlegt vollständig gelöst, ausdrücklich gegliedert bis 75,8\,\% \\
        \bottomrule
    \end{tabularx}
\end{table}

Die letzte Zeile benennt den folgenreichsten Gestaltungshebel. Die zusammengesetzte Aufgabe aus Verknüpfung, Bereinigung und Aggregation löste als Einzelaufgabe kein Modell. Zerlegt und mit fehlerfreier Eingabe lösten dieselben Modelle sie vollständig, ausdrücklich in Teilschritte gegliedert zu bis zu drei Vierteln. Begrenzend wirkt also nicht die Schwierigkeit der einzelnen Operation, sondern die Zahl der gleichzeitig und der unausgesprochen zu treffenden Entscheidungen. Bevor ein Vorhaben ein Modell auswählt oder einen Prompt optimiert, sollte es seine Aufgabenstellung zerlegen \autocite{khanOverviewETLTechniques2024,nagabhyruBridgingTraditionalETL2022}.

\section{Empfehlungen zur Modell- und Prompt-Auswahl}
\label{sec:empfehlungen-auswahl}

Steht fest, dass ein \ac{LLM} eingesetzt werden soll, sind drei Festlegungen zu treffen. Sie werden hier in der Reihenfolge ihrer Wirkung behandelt, die der verbreiteten Aufmerksamkeitsverteilung genau entgegengesetzt ist: Der Ausführungsmodus wirkt am stärksten, die Prompt-Gestaltung deutlich, die Modellwahl am wenigsten.

\noindent\textbf{Ausführungsmodus:} Die Wahl zwischen Code-Generierung und Direktverarbeitung ist keine technische Detailfrage, sondern bestimmt, welche Fähigkeit des Modells überhaupt zur Wirkung kommt. Für den Regelbetrieb ist der Code-Generierungs-Modus vorzuziehen, und zwar aus drei unabhängigen Gründen. Er ist um mehr als eine Größenordnung günstiger, 0,117 bis 0,216~US-Dollar gegenüber 1,14 bis 1,35 für dieselbe Strecke auf der halben Datenmenge. Er erzeugt ein prüfbares, versionierbares Artefakt, das sich einmal abnehmen und danach unverändert ausführen lässt, während die Direktverarbeitung bei jedem Lauf neu und potenziell anders entscheidet. Bei festem Prompt ist sein Modellaufruf weitgehend vom Datenvolumen entkoppelt, Laufzeit und Speicherbedarf des erzeugten Codes hängen dagegen weiterhin von Datenmenge und Algorithmus ab, während der Direkt-Modus durch das Ausgabefenster begrenzt ist und im Rahmen dieser Arbeit eine Halbierung des Datensatzes erzwang. Der Direkt-Modus ist demnach nicht als durchgängiger Verarbeitungsmodus zu empfehlen, sondern gezielt für einzelne, semantisch geprägte Schritte, dort allerdings mit messbarem Gewinn.

\noindent\textbf{Prompt-Gestaltung:} Hier ist die wichtigste Empfehlung eine über die Reihenfolge: Der Zuschnitt der Aufgabe ist vor dem Prompt zu optimieren, nicht danach. Die in Abschnitt~\ref{subsec:gruppen} berichtete Gegenüberstellung zeigt, dass der auffälligste Strategieunterschied der Faktormatrix, 53 gegenüber 72~Prozent bei der Duplikaterkennung, vollständig verschwindet, sobald dieselben Schritte in angemessenem Zuschnitt und mit sauberer Eingabe gestellt werden: Dort liegen alle drei Strategien zwischen 98,8 und 99,4~Prozent. Beispiele und schrittweise Herleitung helfen einem Modell, eine überladene Aufgabe zu bewältigen. Sie werden entbehrlich, sobald die Aufgabe nicht mehr überladen ist. Wer zuerst den Prompt optimiert, arbeitet daher an einem Symptom.

Für verbleibende Teilaufgaben liefern die vorliegenden Messungen Anhaltspunkte zur Auswahl einer Prompt-Vorlage: Bei den unscharfen Duplikaten erreichte die beispielgestützte Vorlage 82,0 gegenüber 26,6~Prozent für Zero-Shot. Bei der schlüsselbasierten Entdoppelung fiel die Vorlage mit expliziter Schrittanleitung günstig aus. Diese Beobachtungen gelten für die konkret eingesetzten Vorlagen und Wiederholungen. Sie ersetzen weder einen Vergleich alternativer Beispiele noch eine Erprobung auf neuen Daten. Wo kein Vorteil beobachtet wird, ist die einfachere Vorlage ein naheliegender Ausgangspunkt.

Zusätzlicher Kontext sollte gezielt ausgewählt und sein Nutzen geprüft werden. Die Angabe der Eingabe-Datentypen behob in der Entwicklungsphase beispielsweise ein stilles Versagen der Spaltenauswahl. Daraus folgt jedoch nicht, dass Kontext ausschließlich nach einem beobachteten Fehler ergänzt werden darf. Auch vorausschauend formulierte Anforderungen können sinnvoll sein, ihr zusätzlicher Nutzen wurde hier nicht separat erhoben. Die Ergebnisse sprechen für eine aufgabenbezogene Auswahl, nicht für eine allgemeine Rangfolge der Prompting-Verfahren.

\noindent\textbf{Modellwahl:} Die drei untersuchten Cloud-Modelle liegen so eng beieinander, dass sich eine allgemeine Empfehlung nicht begründen lässt. Auch die Kosten und Antwortzeiten unterscheiden sich nur moderat. Sinnvoll ist stattdessen eine Auswahl anhand der konkret vorliegenden Aufgabenart, wobei die Modelle in ihren Fehlermustern durchaus unterscheidbar sind. In der gruppierten Strecke, also bei angemessenem Aufgabenzuschnitt, verschwindet auch dieser Unterschied: Dort liegen alle drei zwischen 99,0 und 99,3~Prozent.

 Das lokale Modell wurde nur in einer Konfiguration untersucht: als Llama-Modell mit acht Milliarden Parametern und begrenztem Kontextfenster auf der verfügbaren Hardware. Der beobachtete Rückstand lässt sich deshalb nicht eindeutig in Beiträge von Modell, Parametrisierung und Hardware zerlegen und ist keine Aussage über lokale Modelle insgesamt. Jellyfish zeigt, dass auch die aufgabenspezifische Anpassung lokaler Modelle eine relevante Vergleichsdimension ist \autocite{zhangJellyfish2024}. Für eine praktische Auswahl sollten geeignete lokale Konfigurationen auf den eigenen Aufgaben geprüft und Anschaffungs-, Betriebs- und Prüfkosten mitberücksichtigt werden.

Als Ausgangspunkt einer solchen Erprobung bieten sich klar abgegrenzte Aufgaben und ein vorab festgelegtes Versuchslimit an. Die beobachteten Bestwerte zeigen, dass die untersuchte lokale Konfiguration einzelne Aufgaben vollständig lösen konnte. Sie belegen jedoch weder ein allgemeines Aufgabenverständnis noch, dass größere Hardware allein die übrigen Fehler beheben würde.

\noindent\textbf{Versionsstand:} Die kommerziellen Modelle werden serverseitig gepflegt und lassen sich weder einfrieren noch nach einer Aktualisierung wiederherstellen. Wer eine gemessene Konfiguration produktiv einsetzt, sollte deren Versionskennung protokollieren und die Referenzläufe nach jeder Modellaktualisierung wiederholen.

\noindent\textbf{Messung statt Übernahme:} Keine dieser Empfehlungen ersetzt die Messung im eigenen Kontext, und die Anlage dieser Untersuchung zeigt, welche Anforderungen eine solche Messung erfüllen muss. Eine Konfiguration ist an Daten zu bewerten, gegen die sie nicht justiert wurde (Abschnitt~\ref{sec:versuchsablauf}), und mehrfach auszuführen: Der Endschritt derselben Konfiguration schwankte über fünf identische Wiederholungen zwischen 27,3 und 72,7~Prozent (Abschnitt~\ref{sec:ergebnisse-reproduzierbarkeit}). Wer auf dieser Grundlage eine einzelne Messung zur Entscheidung heranzieht, wählt zu einem erheblichen Teil zufällig.

\section{Hybride Architekturen}
\label{sec:hybride-architekturen}

Aus den vorstehenden Empfehlungen ergibt sich, dass weder ein rein regelbasiertes noch ein rein \ac{KI}-gestütztes Vorgehen der jeweils anderen Variante über alle Aufgaben hinweg überlegen ist. Die naheliegende Konsequenz ist eine Architektur, die beide Verfahren nicht als Alternativen, sondern als Bausteine derselben Strecke behandelt und je Schritt dasjenige einsetzt, dessen Stärken zur Aufgabe passen \autocite{coteDataCleaningMachine2024,kasireddyRoleAIModern2026}.

Der erste Baustein ist die bereits begründete Zerlegung. Sie ist hier nicht nur eine Frage der Übersichtlichkeit, sondern die Voraussetzung dafür, dass je Schritt überhaupt eine Verfahrenswahl möglich wird, und eine mögliche Verbesserung der Bearbeitbarkeit, zusätzlicher Spezifikations- und Prüfaufwand ist dabei zu berücksichtigen.

Der zweite Baustein ist die Zuordnung nach Tabelle~\ref{tab:eignungsmatrix}: regelbasierte Ausführung für die syntaktisch beschreibbaren Schritte, ein Modell im Direkt-Modus für die wenigen semantisch geprägten Schritte, und die abschließende Aggregation wieder regelbasiert oder über generierten Code. Für die untersuchte Strecke bedeutet dies konkret, dass sieben der neun Schritte regelbasiert auszuführen wären und lediglich die Länder-Vereinheitlichung sowie die unscharfe Duplikaterkennung ein Modell benötigen. Eine Beschränkung der Modellaufrufe auf diese Schritte könnte die Modellgebühren reduzieren. Die tatsächlichen Kosten und die Gesamtqualität der hybriden Variante wurden jedoch nicht gemessen.

Der dritte und wichtigste Baustein ist die Prüfung nach jedem Schritt, die unmittelbar aus dem Befund zur Fehlerfortpflanzung folgt. Deren Tücke liegt darin, dass die verursachenden Schritte für sich betrachtet unauffällig wirken: Sie erreichten durchweg über 96~Prozent, und dennoch genügten die wenigen falsch gewählten Duplikat-Repräsentanten, um das Endergebnis um mehr als ein Drittel zu verschieben. Entscheidend ist zudem, dass diese Prüfung nicht auf der Abwesenheit von Fehlermeldungen beruhen darf, wie der Fall des stillen Versagens aus Abschnitt~\ref{sec:diskussion} zeigt. Zu prüfen sind daher inhaltliche Eigenschaften des Ergebnisses, etwa dass die Länderspalte ausschließlich gültige zweistellige ISO-Codes oder den ausdrücklich erlaubten Wert \texttt{UNKNOWN} enthält, dass die Datensatzanzahl im erwarteten Bereich liegt oder dass keine Bestellung ohne zugehörigen Kunden zurückbleibt \autocite{ridzuanReviewDataQuality2024,mumuniAutomatedDataProcessing2025}.

Der vierte Baustein betrifft die kontrollierte Wiederverwendung geprüfter Artefakte. Im Code-Generierungs-Modus kann der geprüfte Code versioniert werden, statt ihn bei jeder Ausführung neu zu erzeugen. Bei einem semantischen Direkt-Schritt können bereits geprüfte Zuordnungen für wiederkehrende Werte gespeichert werden. Neue oder veränderte Eingaben benötigen weiterhin eine Prüfung. Dieses Vorgehen begrenzt wiederholte Generierung und erleichtert die Nachvollziehbarkeit, garantiert aber keine Fehlerfreiheit bei veränderten Daten.

Diese vier Bausteine sind zugleich unter dem Gesichtspunkt des Ressourceneinsatzes zu lesen, wie ihn Abschnitt~\ref{sec:nachhaltigkeit} auswertet. Eine Strecke, die ein Modell nur für zwei von neun Schritten benötigt, verbraucht weniger Rechenleistung, verursacht geringere Nutzungsgebühren und überträgt weniger Daten an externe Anbieter als eine durchgängig modellgestützte. Der vierte Baustein, die Wiederverwendung geprüfter Artefakte, wirkt in dieselbe Richtung, weil er die wiederholte Generierung desselben Codes vermeidet. Nachhaltigkeit und Ergebnisqualität führen bei dieser Architektur folglich zu derselben Gestaltungsentscheidung, was ihre Umsetzung erleichtert: Es ist kein Zielkonflikt aufzulösen \autocite{2026AIIndex,changEfficientPromptingMethods2024}.

Diese Aufteilung beschreibt zugleich die realistische Rolle von \acp{LLM} in der Datenintegration, wie sie sich aus den Ergebnissen dieser Arbeit ergibt. Sie ersetzen die klassische Verarbeitung nicht, sondern verschieben die Grenze dessen, was ohne vorherige vollständige Spezifikation automatisierbar ist. Ihr Beitrag liegt in den wenigen Schritten, in denen eine Vorschrift nicht vollständig formulierbar ist, sowie in der Verkürzung der Zeit, die von der Aufgabenstellung bis zu einer lauffähigen ersten Umsetzung vergeht, und beides ist erheblich, ohne die Aufgabe der prüfenden Instanz zu ersetzen \autocite{walhaDataIntegrationTraditional2024,khanCognitiveETLUsing2025}.
